% Pii — Technical Manual (Developers)
% Template: kaobook
% Compilacao remota via Podman (latex-compiler-shared)
% SSOT-based: architecture + technical analysis live here; operational facts live in the User Manual.

\documentclass[
	fontsize=10pt,
	twoside=false,
	secnumdepth=2,
	open=any,
]{kaobook}

\usepackage[english]{babel}
\usepackage{amsmath, amssymb}
\usepackage{booktabs}
\usepackage{longtable}
\usepackage{array}
\usepackage{listings}
\lstset{
    basicstyle=\ttfamily\scriptsize,
    keywordstyle=\color{blue!70!black},
    commentstyle=\color{green!50!black},
    stringstyle=\color{red!60!black},
    numbers=none,
    breaklines=true,
    frame=none,
    backgroundcolor=\color{blue!5!white},
    showstringspaces=false,
    tabsize=4,
    language=bash
}
\usepackage{tcolorbox}
\tcbuselibrary{skins,breakable}
\newtcolorbox{archbox}[1][]{
    enhanced, breakable,
    colback=blue!5!white, colframe=blue!60!black,
    fonttitle=\bfseries, fontupper=\small, coltitle=white,
    title={#1}, left=5pt, right=5pt, top=5pt, bottom=5pt
}
\newtcolorbox{decbox}[1][]{
    enhanced, breakable,
    colback=green!5!white, colframe=green!60!black,
    fonttitle=\bfseries, fontupper=\small, coltitle=white,
    title={#1}, left=5pt, right=5pt, top=5pt, bottom=5pt
}

\let\refeq\relax
\usepackage{kaotheorems}
\usepackage{kaorefs}
\usepackage{kaobiblio}
\usepackage{xr}
\externaldocument{user-manual}

\graphicspath{{images/}{./}}
\makeindex[columns=3, title=Index, intoc]

\renewcommand*{\chapterformat}{\thechapter\enskip}
\renewcommand*{\chaptermarkformat}{\thechapter\enskip}

\begin{document}

\titlehead{Pii — Technical Manual}
\title[Pii Technical Manual]{%
    Pii\\[0.5em]
    \Large Developer's Manual: Architecture \& Technical Analysis
}
\author[Pii]{Pii Project}
\date{Version 0.1.0}
\publishers{pi-improved}

\frontmatter

\makeatletter
\uppertitleback{\@titlehead}
\lowertitleback{
	\textbf{About this manual} \\\\
	This is the \emph{Technical Manual} for Pii, the single static Rust
	binary that unifies a terminal coding agent, a personal assistant
	daemon, a reactive automation engine, and a self-improving cognitive
	layer.

	\medskip

	\textbf{Audience} \\\\
	Developers who need the architecture decisions and the technical
	analysis.

	\medskip

	\textbf{SSOT} \\\\
	This manual is \emph{SSOT-based}: it holds the architecture and the
	technical analysis. Operational facts (install, configure, run, debug)
	live in the companion \emph{User Manual}. Where a fact has a canonical
	home elsewhere (the \texttt{\_docs/} research tree), this manual links
	to it instead of duplicating it.

	\medskip

	\textbf{Compilation} \\\\
	Compiled remotely via Podman (latex-compiler-shared) on \today.
}
\makeatother

\maketitle

\chapter*{Preface}

Pii is a single static Rust binary. This manual explains \emph{how it
works}: the architecture decisions, the graph engine, the capability and
safety model, the event-sourced journal, the persistence layer, and the
technical analysis that shaped them.

The companion \emph{User Manual} covers installation, configuration,
operation, and debugging for system administrators. The two manuals
complement each other: this one answers \emph{how it works}, the other
answers \emph{how to operate}. Cross-references are marked inline.

\begingroup
\setlength{\textheight}{230\vscale}
\etocstandarddisplaystyle
\etocstandardlines
\tableofcontents
\endgroup

\mainmatter
\setchapterstyle{kao}

%----------------------------------------------------------------------------------------
\chapter{Architecture Overview}
\label{chap:arch}

Pii is a single-process, event-sourced monolith: one redb store,
single-writer; journal + bus + approval gate + budget are coherent
in-process. That monolith is its strength. The full vision and the four
pillars are in the SSOT at
\texttt{\_docs/01-overview/001-pii-vision.md}.

\section{The four pillars}

\begin{enumerate}
    \item \textbf{Coding agent core} (from \texttt{pi\_agent\_rust}) —
        sub-100ms cold start, structured concurrency, capability-gated
        \texttt{Ctx}, JSONL sessions, bash mediation.
    \item \textbf{Assistant/gateway} (from \texttt{zeroclaw}) — channels,
        LLM providers, tools, SOP engine, WASM plugin host, hardware HAL.
    \item \textbf{Cognitive memory + model tiering} (from
        \texttt{yantrikclaw}) — entities/relationships, tier-aware tool
        selection for model sizes 0.5B$\rightarrow$100B+.
    \item \textbf{Self-improvement loop} (from \texttt{hermes-agent}) —
        skill extraction, procedural templates, trajectory logging.
\end{enumerate}

\section{The workspace}

The workspace is a Rust workspace of crates. The kernel ABI is
\texttt{pii-abi} (traits + types only, no transport deps), enforced by
\texttt{xtask deps-check} and a hard file-size gate. The crates:

\begin{longtable}{@{}ll@{}}
\toprule
\textbf{Crate} & \textbf{Role} \\
\midrule
\endhead
\texttt{pii-abi} & kernel ABI: traits + types, no transport deps \\
\texttt{pii-core} & capability context, budget, child-process \\
\texttt{pii-store} & redb persistence (single-writer MVCC ACID) \\
\texttt{pii-graph} & the graph automation engine \\
\texttt{pii-cli} & the \texttt{pii} binary (CLI entry point) \\
\texttt{pii-proto} & UDS JSON-RPC protocol + client \\
\texttt{pii-model} & model providers \\
\texttt{pii-net} & the socket chokepoint / egress policy \\
\texttt{pii-memory} & global memory store + embedder \\
\texttt{pii-worker} & supervised sub-process pool \\
\texttt{pii-tools} & the toolset (fs, edit, bash, glob, grep, http, memory) \\
\texttt{pii-channels} & outbound messaging (Discord, \ldots) \\
\texttt{pii-server} & the daemon: Engine, registry, runs, sessions \\
\texttt{pii-tui} & the terminal UI \\
\texttt{pii-web} & the read-only web viewer \\
\texttt{pii-mcp} & MCP server + client \\
\texttt{pii-uds} & client-side UDS transport helpers (runtime dir, socket connect) \\
\texttt{pii-http-fetch} & credential-free GET-only HTTP fetch tool (via \texttt{pii-net}) \\
\texttt{pii-evolve} & the self-improving boss (separate binary) \\
\texttt{xtask} & build/CI helpers \\
\bottomrule
\end{longtable}

\section{Feature gating}

Pii is intentionally a monolith, but the installer may want a smaller or
more hardened binary. Cargo \textbf{features} are the lever for
\emph{what's in the build} — NOT for \emph{what's a microservice}. A
feature trims the compile; it does not create a process/IPC boundary. The
convention: gate surfaces with risk or optional deps; never gate the
engine substrate. See \texttt{\_docs/02-guides/feature-gating.md}.

\begin{decbox}[Design decision]
Two subsystems are a security win as opt-in features: \texttt{channels},
\texttt{tls}, and \texttt{github} all do network egress and may carry a
token. A feature that \emph{excludes the code} is stronger than a runtime
config that \emph{disables} it — the dialing code isn't compiled in, so it
cannot be reached.
\end{decbox}

%----------------------------------------------------------------------------------------
\chapter{The Capability Context}
\label{chap:ctx}

The enforcing capability context is \texttt{pii-core::Ctx}. It combines
four concerns:

\begin{itemize}
    \item \textbf{Caps} — \texttt{CapSet} with monotonic attenuation via
        intersection.
    \item \textbf{Budget} — carveable, never-duplicated accounting.
    \item \textbf{Cancel} — an ancestor-propagating \texttt{CancelToken}.
    \item \textbf{Ctx} — the enforcing context that combines all four.
\end{itemize}

\section{Capability bits}

The coarse capability bits live in \texttt{pii-abi} (so the capability
surface is part of the stable ABI — a WASM/worker node or plugin can
reference \texttt{CapKind} without depending on \texttt{pii-core}):

\begin{longtable}{@{}lll@{}}
\toprule
\textbf{Bit} & \textbf{Value} & \textbf{Meaning} \\
\midrule
\endhead
\texttt{FS\_READ} & 1 & read files within an allow path scope \\
\texttt{FS\_WRITE} & 2 & write files within an allow path scope \\
\texttt{NET} & 4 & reach the network \\
\texttt{PROC} & 8 & spawn/talk to processes \\
\texttt{HW} & 16 & touch hardware \\
\texttt{SECRETS} & 32 & access secrets \\
\texttt{MEMORY\_WRITE} & 64 & write into the agent's memory \\
\texttt{SPAWN} & 128 & spawn subprocesses \\
\bottomrule
\end{longtable}

Values must remain stable forever.

\section{The \texttt{Ctx} accessors}

\begin{longtable}{@{}ll@{}}
\toprule
\textbf{Method} & \textbf{Purpose} \\
\midrule
\endhead
\texttt{root(caps, budget)} & build a root context \\
\texttt{fs()} & filesystem capability (scoped) \\
\texttt{net()} & network capability (scoped) \\
\texttt{proc()} & process capability (scoped) \\
\texttt{child(principal, narrow, share)} & monotonically-attenuating child \\
\texttt{child\_with\_token(...)} & child with a cancel token \\
\texttt{cancelled()} & is the run cancelled? \\
\texttt{elapsed\_ms()} & elapsed since start \\
\texttt{now\_ms()} & current clock time \\
\texttt{deadline\_ms()} & the deadline, if any \\
\texttt{check()} & cancellation + deadline seam \\
\texttt{observer()} & the observer \\
\texttt{caps()} & the capability set \\
\texttt{budget()} & the budget \\
\texttt{principal\_str()} & the principal \\
\texttt{principal()} & the principal \\
\texttt{trace()} & the trace context \\
\texttt{cancel\_token()} & the cancel token \\
\bottomrule
\end{longtable}

\section{The budget}

The budget is carveable and never-duplicated, across four dimensions:
\texttt{Tokens}, \texttt{ToolCalls}, \texttt{MicroUsd}, \texttt{WallMs}.
A \texttt{BudgetShare} is \texttt{All}, \texttt{Fraction(pct)},
\texttt{Absolute\{...\}}, or \texttt{Shared} (same-process nodes share the
parent's budget without carving). \texttt{charge()} returns
\texttt{BudgetExhausted} when the amount exceeds what remains, leaving the
counter at 0.

%----------------------------------------------------------------------------------------
\chapter{The Graph Engine}
\label{chap:graph}

Pii ships a deterministic, statically-validated \textbf{graph automation
engine} (\texttt{crates/pii-graph}) that is the control plane under every
agent behavior. A graph is a dataflow/control-flow program expressed in
TOML whose execution is \textbf{event-sourced}: every step appends to an
append-only run journal, and a run can be resumed by replaying that
journal.

A graph is NOT a script. It is a typed assembly of \textbf{nodes} (pure
functions, LLM agents, tools, approval gates, timers), connected by
\textbf{edges} that carry data or ordering, with explicit \textbf{join
policies} for fan-in, \textbf{route} nodes for if/then/else,
\textbf{loop budgets} for bounded cycles, \textbf{capability declarations}
at three layers, and \textbf{effects declarations} that decide parallel vs
serial execution.

The full engine design is in the SSOT at
\texttt{\_docs/06-research/006-pii-graph-engine.md}. The operational
commands (\texttt{pii graph lint|show|run|resume|trace}) are in the User
Manual, Chapter~\ref{chap:run}.

\section{The specification model}

\begin{lstlisting}
GraphSpec { name, entry: NodeId, caps, limits, nodes: [NodeSpec], edges: [Edge] }
NodeSpec  { id, kind, join: JoinPolicy, loop_budget: Option<u32>,
            retry: RetryPolicy, caps, cfg: serde_json::Value }
Edge      { from, port: PortName, to, kind: EdgeKind }
\end{lstlisting}

\begin{itemize}
    \item \textbf{Edges}: \texttt{EdgeKind::Data} (value-carrying
        dependency: readiness + payload) vs \texttt{EdgeKind::Order}
        (ordering-only: readiness, no payload).
    \item \textbf{JoinPolicy}: \texttt{All} (default — every inbound edge
        must fire), \texttt{Any \{ cancel\_losers \}} (any one inbound edge
        fires; optionally cancels sibling branches), \texttt{Quorum \{ k,
        cancel\_losers \}} (at least \texttt{k} fire), and
        \texttt{AllSettled} (wait for every inbound edge to settle, fired
        or dead). The compact string grammar: \texttt{"all"},
        \texttt{"any"}, \texttt{"any:cancel"}, \texttt{"quorum:K"},
        \texttt{"quorum:K:cancel"}, \texttt{"allsettled"}.
    \item \textbf{Limits}: \texttt{max\_steps} (default 1000),
        \texttt{max\_wall\_ms} (default 600000),
        \texttt{max\_tokens} (default unbounded),
        \texttt{max\_micro\_usd} (default unbounded),
        \texttt{max\_tool\_calls} (default unbounded),
        \texttt{max\_parallel} (default 0 = available parallelism).
        Exhaustion is terminal and non-retryable.
    \item \textbf{RetryPolicy}: per-node \texttt{max} retries on
        retryable failures.
    \item \textbf{Expression language}: deliberately tiny —
        \texttt{\$.path.into.context} object-member access only; no
        arrays, no functions. Used by \texttt{transform.expr},
        \texttt{route.on}, and tool args.
\end{itemize}

\section{Node kinds}

\begin{longtable}{@{}llll@{}}
\toprule
\textbf{Kind} & \textbf{Layer} & \textbf{Role} & \textbf{Output ports} \\
\midrule
\endhead
\texttt{transform} & pure built-in & JSON reshape via \texttt{expr} & ok / err \\
\texttt{route} & pure built-in & branch on \texttt{on}-expr via \texttt{rules[]} & one per rule \\
\texttt{emit} & built-in & write a value to a file (fs\_write cap) & ok / err \\
\texttt{wait} & pure built-in & sleep/fire-or-timeout & fired / timeout \\
\texttt{timer} & pure built-in & delay \texttt{in\_ms} or 5-field \texttt{cron} & ok \\
\texttt{agent} & server & LLM agent: system prompt, tools allowlist, \texttt{max\_turns}, \texttt{loop\_budget}, compaction & ok / err / budget \\
\texttt{tool} & server & human-authored command (e.g.\ \texttt{bash}) & ok / err \\
\texttt{gate} & server & suspend the run for human approval & ok \\
\texttt{heal\_detect} & server & self-healing ledger trigger & ok \\
\texttt{memory} & server & standalone \texttt{consolidate} maintenance op & ok \\
\texttt{worker} & server & supervised sub-process pool & ok \\
\texttt{spawn} & server & launch a child run (spec 32) & ok / escalate \\
\texttt{crash} & debug & deliberate failure for tests & --- \\
\bottomrule
\end{longtable}

\section{If/Then/Else — the \texttt{route} node}

Branching is not a separate ``cond'' node: \texttt{route} maps a
\texttt{\$.} expression to output ports via an ordered rule list, and
edges wire each port to a different downstream node. The last rule
without an \texttt{eq} is the else.

\begin{lstlisting}
[[node]]
id = "gate"
kind = "route"
on = "$.payload.exit_code"
rules = [
  { eq = 0, port = "pass" },
  { port = "fail" },          # fallthrough / else
]
\end{lstlisting}

\section{Loops — bounded cycles via \texttt{loop\_budget}}

Cycles are legal when a member node declares \texttt{loop\_budget = N}.
\texttt{is\_in\_cycle} (tarjan SCC) detects cycles; \texttt{reset\_loop}
re-runs the loop body genuinely instead of reading stale state. Exceeding
the budget is a terminal failure.

\begin{lstlisting}
fix [agent] (entry) (loop_budget=3)
 |-- ok --> test
test [tool: bash]
 |-- ok --> gate
gate [route: $.payload.exit_code]
 |-- pass --> done
 |-- fail --> retry
retry [transform: $.payload]
 ^-- loops back to "fix" (budget=3)
done [transform: $.payload]
\end{lstlisting}

\section{Capabilities and Effects (safety model)}

\textbf{Capabilities} (\texttt{crates/pii-core/src/cap.rs}): bit flags —
FS\_READ, FS\_WRITE, PROC, NET, \ldots Effective node permissions are a
\textbf{three-layer intersection}: \texttt{run caps $\cap$ graph caps
$\cap$ node caps}. A denial at any layer fails the node with
\texttt{capability denied}. This is a real bug class: in 2026-09-05 the
live \texttt{ollama:} provider failed because NET was missing at the run
layer AND at the graph/node layers — all three had to be fixed for a real
LLM call to work. (Operational symptom: see the User Manual,
Chapter~\ref{chap:debug}.)

\textbf{Effects} (\texttt{crates/pii-abi/src/effects.rs}): each node kind
declares static effect classes (PURE, READS\_FS, WRITES\_FS, NET,
SPAWNS\_PROC, MUTATES\_EXTERNAL, LONG\_RUNNING). The scheduler partitions
each batch: \texttt{PURE | READS\_FS} run in parallel (bounded by
\texttt{max\_parallel}), \texttt{WRITES\_FS | MUTATES\_EXTERNAL} run
sequentially. Effects fail closed — a node that does not declare its
effects is treated as fully serialising.

\section{Execution: tick loop, batch drain, event-sourced journal}

\begin{enumerate}
    \item \textbf{Tick loop}: propagate dead branches to a fixpoint
        (journaling each new \texttt{NodeDead}); fail LOUDLY on an
        unsatisfiable join; drain ALL ready nodes into one batch; check
        loop budgets; dispatch.
    \item \textbf{Partition}: parallel-safe nodes run in bounded thread
        chunks; serialising nodes run in order. Results apply as they
        arrive; \texttt{cancel\_losers} runs after each result.
    \item \textbf{Journal}: EVERY event appends to the run's append-only
        journal in redb — \texttt{Started}, \texttt{TickStarted},
        \texttt{NodeStarted} (journaled BEFORE the node runs, so an
        interrupted node is visible as ``may have executed''),
        \texttt{NodeFinished}, \texttt{NodeFailed}, \texttt{NodeCancelled},
        \texttt{NodeDead}, \texttt{NodeInterrupted} (audit marker on
        resume), \texttt{Suspended}, \texttt{Resumed}, \texttt{Finished},
        \texttt{RunSpawned}. A header (status + next-seq) is persisted
        with each append; the blackboard snapshots every 32 events and
        unconditionally on suspend.
    \item \textbf{Terminal dispositions}: Running $\rightarrow$ Succeeded
        | Failed | Suspended (human approval) $\rightarrow$ Resumed.
    \item \textbf{Resume}: reads the journal, recovers the input from the
        content-addressed store, replays every event to rebuild state,
        re-propagates dead, appends \texttt{Resumed}, and continues. A
        node with a terminal event is never re-invoked; a
        \texttt{NodeStarted} with no terminal event re-runs.
\end{enumerate}

The \textbf{seq} counter is the single source of run ordering — the same
number the TUI shows as the live \texttt{[seq 0x\ldots]} indicator and
that \texttt{pii graph trace} prints per event. (Operational: see the User
Manual, Chapter~\ref{chap:debug}.)

%----------------------------------------------------------------------------------------
\chapter{Persistence}
\label{chap:store}

Pii uses \textbf{redb} (single-writer MVCC ACID, pure Rust, no C dep for
static musl). Consequence: the daemon owns the DB; a UDS control channel
is the access path. The store is a single-file redb database.

The daemon is the single writer. The control socket is local-only Unix
(mode 0600), and peer-credential authorization rejects a peer whose uid
differs from the daemon's. (Operational: see the User Manual,
Chapter~\ref{chap:ops}.)

\section{The single-writer rule}

A \texttt{redb::WriteTransaction} is opened and committed inside one
method body. It is never stored in a struct, never returned, and never
held across a node dispatch. This invariant is what makes the journal's
``one write txn per event'' property hold.

\section{The append-only journal}

The run journal is an append-only log of typed \texttt{RunEvent}s — the
single source of truth for the whole run. The LLM message history is
derived by pure projection, never stored separately; replay is
re-derivation. This is the same pattern DeepSeek Harness (dsh) uses, and
it is the emergent standard for agent state management — two independent
implementations (dsh, pidag/pi) validate it. See
\texttt{\_docs/06-research/005-dsh-vs-pii-harness-analysis.md}.

%----------------------------------------------------------------------------------------
\chapter{The Control Plane}
\label{chap:control}

\texttt{pii-server} (\texttt{Engine}) owns the registry (built-ins +
agent/tool/gate/heal/memory/worker/spawn), spawns runs
(\texttt{run.start}, \texttt{run.resume}, \texttt{run.cancel}), tracks
sessions, and streams events (\texttt{event.subscribe} — replay-since +
live, bounded queue, \texttt{Lagged} on overflow).

The daemon socket is local-only Unix (0600). The full method surface:

\begin{longtable}{@{}ll@{}}
\toprule
\textbf{Method} & \textbf{Purpose} \\
\midrule
\endhead
\texttt{sys.hello} & handshake \\
\texttt{sys.status} & server status \\
\texttt{sys.version} & server version \\
\texttt{sys.shutdown} & shut the daemon down \\
\texttt{sys.reload} & reload config (channels/workers), never restarts \\
\texttt{session.create} & create a session \\
\texttt{session.prompt} & prompt a session \\
\texttt{session.interrupt} & interrupt the session's active run \\
\texttt{session.state} & session snapshot \\
\texttt{session.list} & list sessions \\
\texttt{session.close} & close a session \\
\texttt{run.start} & start a graph run \\
\texttt{run.get} & run header \\
\texttt{run.list} & list runs \\
\texttt{run.cancel} & cancel a run \\
\texttt{run.resume} & resume an interrupted run \\
\texttt{run.trace} & fetch a run's journal trace \\
\texttt{graph.apply} & store/replace a graph \\
\texttt{graph.get} & fetch a graph \\
\texttt{graph.list} & list graphs \\
\texttt{graph.lint} & lint a graph source \\
\texttt{event.subscribe} & subscribe to events \\
\texttt{event.unsubscribe} & unsubscribe \\
\texttt{worker.list} & list workers and their state \\
\texttt{worker.reload} & reload a worker (drain $\rightarrow$ respawn) \\
\texttt{approve.list} & list pending approvals \\
\texttt{approve.resolve} & resolve an approval (allow/deny) \\
\texttt{heal.list} & list heal-attempt ledger entries \\
\texttt{heal.get} & fetch a heal entry \\
\texttt{cron.list} & list cron entries \\
\texttt{channel.list} & list configured channels \\
\texttt{memory.search} & memory search \\
\texttt{memory.put} & memory put \\
\texttt{memory.consolidate} & memory consolidate \\
\bottomrule
\end{longtable}

The thin CLI subcommands that call these methods are in the User Manual,
Chapter~\ref{chap:run}.

\section{The wire protocol}

The wire contract is \texttt{pii-proto}: newline-delimited JSON (JSON-RPC
2.0 shape), protocol version \texttt{pii.v2}. Frames are
\texttt{Request}, \texttt{Response}, or \texttt{Notification}. Error
codes: JSON-RPC reserved (\texttt{-32700} ParseError, \texttt{-32600}
InvalidRequest, \texttt{-32601} MethodNotFound, \texttt{-32602}
InvalidParams, \texttt{-32603} Internal) plus Pii's application range
(\texttt{100} ProtocolMismatch, \texttt{101} Unauthorized, \texttt{102}
NotFound, \texttt{103} Denied, \texttt{104} BudgetBreach, \texttt{105}
Cancelled, \texttt{106} Busy).

\section{The push vocabulary}

The server$\rightarrow$client push vocabulary (\texttt{pii-proto::event})
is a closed, \texttt{kind}-tagged enum: \texttt{RunStarted},
\texttt{NodeStarted}, \texttt{NodeFinished}, \texttt{NodeFailed},
\texttt{NodeDead}, \texttt{AssistantDelta}, \texttt{ToolCall},
\texttt{ToolResult}, \texttt{RunFinished}, \texttt{RunCancelled},
\texttt{RunSuspended}, \texttt{RunResumed}, \texttt{ApprovalRequested},
\texttt{ApprovalResolved}, \texttt{WorkerDied}, \texttt{WorkerRestarted},
\texttt{ChannelMessage}, \texttt{MemoryDegraded}, \texttt{RouteDecision},
\texttt{RunSpawned}, \texttt{Lagged}. A subscriber that falls behind gets
a \texttt{Lagged} event (never silent).

%----------------------------------------------------------------------------------------
\chapter{The Model Layer}
\label{chap:model}

\texttt{pii-model} is the model provider seam. \texttt{ModelProvider} is
\textbf{sync} on purpose: \texttt{GraphNode::run} is sync and runs execute
on a blocking thread. Providers never read configuration themselves except
from the explicit spec string handed to \texttt{provider\_from\_env} by
the CLI.

\section{Chat types}

\begin{itemize}
    \item \texttt{Role}: System, User, Assistant, Tool.
    \item \texttt{ToolCall}: id, name, JSON arguments.
    \item \texttt{ChatMessage}: role, content, optional tool\_call\_id,
        tool\_calls.
\end{itemize}

\section{Providers}

The shipped provider is \texttt{OllamaProvider} (via \texttt{ollama.rs}).
The provider spec string (e.g.\ \texttt{"ollama:granite4.2:3b"}) selects
the provider and model. \texttt{provider\_from\_env} builds a provider
from the environment.

\section{Model routing (spec 28)}

The \texttt{agent} node can declare \texttt{route\_step}; resolution goes
through \texttt{resolve\_model}, which looks the step up in
\texttt{[evolve.routing]}. A \texttt{route\_step} with no entry is a
\textbf{loud} \texttt{NotFound} — never a silent default. The known steps
are \texttt{localize}, \texttt{patch}, \texttt{review}, \texttt{research},
\texttt{orchestrator}. Each routing decision is broadcast as a
\texttt{RouteDecision} event on the live bus (the journal is the audit
trail).

%----------------------------------------------------------------------------------------
\chapter{The Toolset}
\label{chap:tools}

\texttt{pii-tools} is the toolset. The \texttt{Tool} trait has
\texttt{name()}, \texttt{spec()} (model-facing), and \texttt{call()}
(execution) — the type split makes leakage unrepresentable. The built-in
set:

\begin{longtable}{@{}ll@{}}
\toprule
\textbf{Tool} & \textbf{Purpose} \\
\midrule
\endhead
\texttt{fs\_read} & read a file within the allow scope \\
\texttt{fs\_write} & write a file within the allow scope \\
\texttt{edit} & edit a file \\
\texttt{multi\_edit} & multiple edits in one call \\
\texttt{grep} & search file contents \\
\texttt{glob} & find files by name \\
\texttt{bash} & run a command (via \texttt{/bin/sh}, output capped at 64 KiB) \\
\texttt{worktree} & provision a disposable git worktree (spec 32) \\
\bottomrule
\end{longtable}

The \texttt{bash} tool caps output at \texttt{MAX\_OUTPUT = 64 KiB} and
runs only inside the workspace. All process execution goes through
\texttt{Ctx::proc()}.

%----------------------------------------------------------------------------------------
\chapter{Memory}
\label{chap:memory}

\texttt{pii-memory} is the global memory store: BM25 + i8 vectors + graph
lane fused by RRF. Zero model calls, zero socket code; redb is reached
only through \texttt{pii\_store::MemStore}.

\section{Layers and kinds}

\begin{itemize}
    \item \textbf{Layers}: L1 (atoms — one fact with evidence), L2 (scenes
        — distilled summaries of related atoms), L3 (skills and persona).
    \item \textbf{Kinds}: Atom, Scene, Persona, Skill.
\end{itemize}

A \texttt{Doc} has a blake3-content id (idempotent put: the same text put
twice IS the same doc), layer, kind, text, concept tags, importance,
created/last-used times, optional evidence \texttt{(run\_id, seq)}, and an
archived flag.

\section{Search}

\texttt{MemorySearcher::search} fuses BM25, the i8 vector lane, and the
graph lane by RRF. The embedder is \texttt{TrigramEmbedder} (offline,
dims 256, quantised to i8 + scale) by default; \texttt{NullEmbedder} and
\texttt{OllamaEmbedder} also exist. A degraded vector lane surfaces a
\texttt{MemoryDegraded} event (never silent).

%----------------------------------------------------------------------------------------
\chapter{Workers}
\label{chap:worker}

\texttt{pii-worker} is supervision and hot reload. The worker wire
vocabulary lives here, decoupled from \texttt{pii-proto}.

\section{Worker spec}

A \texttt{WorkerSpec} has: name, program, args, env, max\_in\_flight
(default 16), ready\_timeout\_ms (default 2000), restart\_backoff\_ms
(default 50), restart\_window\_ms (default 60000), and
max\_restarts\_in\_window (default 5).

\section{Lifecycle}

\texttt{WorkerState}: Starting (waiting for \texttt{hello}), Ready
(serving invokes), Draining (for a hot reload; new invokes queued),
Restarting (after death), Failed (crash-loop cap hit; invokes refused).

\section{Hot reload}

\texttt{worker.reload} drains the worker, swaps the child (the program
replaced on disk = hot reload), and drains queued requests after the swap.
The \texttt{ReloadReport} carries old/new pid, drained count, and
in-flight at start.

%----------------------------------------------------------------------------------------
\chapter{Channels}
\label{chap:channels}

\texttt{pii-channels} is inbound polling + outbound reply, no sockets.
Every channel goes through \texttt{Arc<dyn HttpTransport>} (from
\texttt{pii-core}), which keeps the crate testable with a fake and keeps
\texttt{cargo xtask net-check} clean.

\section{The \texttt{Channel} trait}

\begin{itemize}
    \item \texttt{name()} — the configured channel name.
    \item \texttt{kind()} — \texttt{"discord"}, \texttt{"cli"}, \ldots
    \item \texttt{poll(ctx, cursor)} — fetch messages newer than a cursor.
    \item \texttt{send(ctx, to, text)} — send a reply.
\end{itemize}

\section{Discord}

\texttt{DiscordChannel} is built from the environment
(\texttt{PII\_DISCORD\_TOKEN}, \texttt{PII\_DISCORD\_CHANNEL\_ID},
\texttt{PII\_DISCORD\_API\_BASE}).

\section{CLI channel}

\texttt{CliChannel} is newline-delimited JSON on two files (inbox/outbox),
no I/O beyond them. Used by tests and \texttt{pii -p} style scripting.

%----------------------------------------------------------------------------------------
\chapter{The Self-Improving Boss}
\label{chap:evolve}

The self-improving boss (\texttt{pii-evolve}) is a separate-binary crate
that runs the boss loop as a real pii graph on the daemon — dogfooding
(012 \S3/\S4). Every 30 minutes it runs the silent-noop-check gate on the
daemon itself; a detected failure escalates through a \texttt{gate} node
so a human approves before anything heals. No FS\_WRITE anywhere in the
watch graph; no agent/model node — detection, forwarding, and the
approval decision are all deterministic nodes.

The deep design is in the SSOT at
\texttt{\_docs/06-research/012-p47-remaining-opus-plan.md}. The opus review and the
remaining-work plan are merged into the same document (\texttt{012-}).

\section{The launcher}

The boss's outer loop is a thin UDS launcher (\texttt{driver.rs}): it
loads \texttt{state.md}, bumps \texttt{cycle\_count}, sends
\texttt{run.start} for \texttt{starter/boss-watch} over the daemon's
control socket, records the launched run id + \texttt{phase =
"watching"} back into \texttt{state.md}, and returns. It spawns no
subprocess and shells out to nothing — the daemon owns detection, gating,
and the approval ledger.

\section{The state doc}

\texttt{state.md} is the boss's durable ``where am I'' source of truth:
goal, phase, last\_step, next\_step, a small properties map, and
updated\_at. Written at every milestone; never inferred from transcript
history. Durable facts (lessons, resolutions, model-routing learnings)
belong in external memory, not here.

\section{The heal-fix loop}

The heal-fix graph is a complete localize$\rightarrow$patch$\rightarrow$
validate$\rightarrow$escalate loop: \texttt{probe} (bash, runs the test
command in a worktree) $\rightarrow$ \texttt{route\_probe} on exit code
$\rightarrow$ \texttt{fix} (agent with fs tools, no bash)
$\rightarrow$ \texttt{retest} $\rightarrow$ \texttt{route\_retest}
$\rightarrow$ \texttt{summarize} (git diff) $\rightarrow$ \texttt{escalate}
(gate) $\rightarrow$ \texttt{done}, with a \texttt{retry} back-edge and
\texttt{loop\_budget = 2}.

\section{Protected paths}

The boss cannot patch its own oracle. \texttt{PathScope} gains a
deny-wins builder; \texttt{run\_caps} adds a protected-path deny list for
scoped runs; \texttt{heal-collect} refuses test-deletion/protected-path
patches. This is the ordering guardrail that prevents a model from
deleting the test that proves it worked.

%----------------------------------------------------------------------------------------
\chapter{Technical Analysis}
\label{chap:analysis}

\section{Convergent architecture (dsh comparison)}

DeepSeek Harness (dsh) validates pi-improved's core architecture
decisions — event-sourced journal, derived state, capability attenuation,
append-only durability — done in TypeScript-with-plugins; pi-improved is
the same pattern in a static Rust binary with a real graph engine on top.
The transferable upgrades are the engineering details dsh has refined
further (write-behind batching, reader-side interrupted-turn repair,
per-tool timeouts, single-writer leases) — not its plugin metaphysics.
See \texttt{\_docs/06-research/005-dsh-vs-pii-harness-analysis.md}.

\section{Context management (Slice C)}

Context management is the one P4.7 slice that is unbuilt. The reference
design to steal is \texttt{dsh-context} (a DeepSeek Harness plugin):
context insight should be a \textbf{projection over the existing journal},
reading the harness's own token accounting, with a delta-mode diagnostic
that directly targets the \texttt{compaction.rs:69} no-op. See
\texttt{\_docs/06-research/013-dsh-context-slice-c-reference.md}.

\section{The graph engine vs. git log}

The lane-allocation algorithm from \texttt{git log --graph} was ported to
render pii graphs as ASCII. The one thing that does not port: loops. Git's
DAG has no cycles; pii graphs genuinely cycle (\texttt{loop\_budget}).
Design decision: lay out nodes in topological order and render each loop
closure as a separate annotation line, never as a lane. See
\texttt{\_docs/06-research/006-pii-graph-engine.md}.

\section{Where to go for more}

\begin{itemize}
    \item \textbf{Operation}: the companion User Manual.
    \item \textbf{Research SSOT}: \texttt{\_docs/06-research/} (numbered
        research write-ups).
    \item \textbf{Guides}: \texttt{\_docs/02-guides/}.
    \item \textbf{Vision}: \texttt{\_docs/01-overview/001-pii-vision.md}.
\end{itemize}

\backmatter
\setchapterstyle{plain}

\printindex

\end{document}
